% ============================================================================
% ALPHA_MODEL_BOOKLET.tex — Alpha Model Reference (Condensed)
% VSEPR-SIM 3.0.0
% ============================================================================
\documentclass[11pt,a4paper]{article}
\usepackage[utf8]{inputenc}
\usepackage[T1]{fontenc}
\usepackage{amsmath,amssymb}
\usepackage{booktabs,longtable}
\usepackage{geometry}
\usepackage{hyperref}
\geometry{margin=2.4cm}

\title{Alpha Model Reference\\{\large VSEPR-SIM 3.0.0 --- Condensed Booklet}}
\author{VSEPR-SIM Project}
\date{\today}

\begin{document}
\maketitle

% ====================================================================
\section{Overview and Production Architecture}

The Alpha Model predicts scalar isotropic polarizability
$\alpha$ (\AA${}^3$) for any element by atomic number $Z$.
Production code lives in
\texttt{atomistic/models/alpha\_model.hpp}.

\medskip
\noindent\textbf{Dual-path architecture:}

\begin{center}
\begin{tabular}{lll}
\toprule
\textbf{Path} & \textbf{Range} & \textbf{Error} \\
\midrule
Empirical table  & $Z=1$--118 & 0\% (by construction) \\
Generative model & $Z>118$    & \(\sim\)15--17\% global RMS (historical best) \\
\bottomrule
\end{tabular}
\end{center}

For $Z \le 118$, $\alpha$ is returned directly from
\texttt{ALPHA\_EMPIRICAL[Z-1]}, a 118-element compile-time table.
For $Z > 118$, the generative fallback fires.

% ====================================================================
\section{Empirical Table}

The table contains experimentally determined and high-level ab initio
polarizabilities for every known element:

\begin{itemize}
  \item \textbf{Z = 1--54}: gas-phase experimental values from
        Miller, \textit{JACS} \textbf{112}, 8533 (1990).
  \item \textbf{Z = 21--118}: DHF/CCSD(T) recommended values from
        Schwerdtfeger \& Nagle, \textit{Mol.\ Phys.}\ \textbf{117}, 1200 (2019).
\end{itemize}

For overlapping elements (Z = 21--54), experimental values are preferred
where available; computational values fill the heavy-element range.
The table is the sole prediction path in production and carries 0\%
model error by construction.

% ====================================================================
\section{Generative Fallback Model}

The generative model exists only for hypothetical super-heavy elements
($Z > 118$).  During development it was also trained against the
Z = 1--118 reference data to validate its functional form.
It never exceeded 15\% global RMS accuracy on that training set.

\subsection{Model Equation}

\begin{equation}
\boxed{\;
  \alpha(Z) \;=\; r_\mathrm{eff}(Z)^{\,3}\;\cdot\;
  g_\mathrm{block}(Z)\;\cdot\;
  g_\mathrm{period}(Z)\;\cdot\;
  g_f\!\bigl(n_f(Z)\bigr)\;\cdot\;
  g_\mathrm{bind}(Z)
  \;+\; \Delta_f(Z)
\;}
\label{eq:alpha}
\end{equation}

\noindent where each gate factor is:

\begin{align}
  r_\mathrm{eff} &= r_\mathrm{cov}\,(1 + c_\mathrm{rel}\,Z^2)\,(1 - \beta_f\,n_f/14)
  \label{eq:reff} \\[4pt]
  g_\mathrm{block} &= c_{\mathrm{block}}[\,b(Z)\,], \quad b \in \{s,p,d,f\}=\{0,1,2,3\}
  \\[4pt]
  g_\mathrm{period} &= k_{\mathrm{period}}[\,p(Z)-1\,], \quad p \in 1\ldots7
  \\[4pt]
  g_f(n_f) &= 1 + a_{f1}\frac{n_f}{14}
             + a_{f2}\exp\!\Bigl[-\Bigl(\frac{n_f-7}{\sigma_{f1}}\Bigr)^{\!2}\Bigr]
             + a_{f3}\exp\!\Bigl[-\Bigl(\frac{n_f-14}{\sigma_{f2}}\Bigr)^{\!2}\Bigr]
  \label{eq:gf} \\[4pt]
  g_\mathrm{bind} &= \frac{1}{1 + b_\mathrm{bind}\,I_1(Z)}
  \label{eq:gbind}
\end{align}

The additive blob correction $\Delta_f$ fires only for f-block elements:
\begin{equation}
  \Delta_f(Z) = b_{f,\mathrm{lin}}\,n_f
              + b_{f,\mathrm{half}}\,\delta_\mathrm{half}(Z)
              + b_{f,\mathrm{full}}\,\delta_\mathrm{full}(Z)
  \label{eq:blob}
\end{equation}

\subsection{Physical Interpretation of Gate Factors}

\begin{description}
  \item[$r_\mathrm{eff}$ (Eq.~\ref{eq:reff})]
    Covalent radius with relativistic contraction ($c_\mathrm{rel} < 0$)
    and lanthanide/actinide contraction ($\beta_f$): 4f/5f electrons are
    poor shielders, so $Z_\mathrm{eff}$ rises with filling.
  \item[$g_f$ (Eq.~\ref{eq:gf})]
    Shell-closure gate.  Two Gaussians at half-shell ($n_f=7$: Eu, Am)
    and full-shell ($n_f=14$: Yb, No) capture subshell-stabilization
    anomalies.
  \item[$g_\mathrm{bind}$ (Eq.~\ref{eq:gbind})]
    First ionization energy as a stiffness measure: tightly bound
    electrons polarize less.
  \item[$\Delta_f$ (Eq.~\ref{eq:blob})]
    Additive correction for electronic-configuration discontinuities that
    the smooth $r_\mathrm{cov}^3$ basis cannot represent.
\end{description}

\subsection{Fitted Parameters (v2.6.2 snapshot)}

15 free parameters in total (7 period + 4 block + 4 model).

\begin{center}
\begin{tabular}{llr}
\toprule
\textbf{Symbol} & \textbf{Role} & \textbf{Value} \\
\midrule
$k_{\mathrm{period},1\text{--}7}$ & period scaling & 20.82, 5.38, 5.12, 3.96, 3.73, 2.92, 2.67 \\
$c_{\mathrm{block},s/p/d/f}$    & block correction & 1.46, 0.97, 1.61, 1.75 \\
$c_\mathrm{rel}$                & relativistic contraction & 0.0 (cold) \\
$\beta_f$                        & f-shielding contraction & 0.0 (cold) \\
$b_\mathrm{bind}$               & ionization stiffness & 0.0 (cold) \\
$q_\mathrm{drude}$              & Drude coupling charge & 1.0\,\textit{e} \\
\bottomrule
\end{tabular}
\end{center}

\noindent Parameters marked `cold'' were never moved from their initial
values during fitting---the optimizer found no gradient signal through them
at the resolution available.

% ====================================================================
\section{Performance Summary}

\subsection{Best Achieved Accuracy}

Some element blocks approached good tolerance while others remained
stubbornly above the 15\% floor.

\begin{center}
\begin{tabular}{lr}
\toprule
\textbf{Block / Subset}  & \textbf{Best RMS (\%)} \\
\midrule
Halogens (F, Cl, Br, I, At)       & 14.2 \\
Period 6 (Cs--Rn)                  & 14.1 \\
s-block (alkali + alkaline earth)  & 14.8 \\
p-block (main-group non-metals)    & 15.0 \\
d-block (transition metals)        & 15.5 \\
f-block (lanthanides + actinides)  & 17.3 \\
\midrule
\textbf{Global}                    & \textbf{15.3} \\
\bottomrule
\end{tabular}
\end{center}

\subsection{Version Progression}

\begin{center}
\begin{tabular}{lrl}
\toprule
\textbf{Version} & \textbf{Global RMS} & \textbf{Key Change} \\
\midrule
v2.0 & 34\% & $r_\mathrm{cov}^3$ only \\
v2.2 & 27\% & + period/block gates \\
v2.4 & 21\% & + ionization binding \\
v2.6 & 17\% & + f-block Gaussians \\
v2.6.2 & 15.3\% & + additive blob \\
v2.8 & 15.1\% & + relativistic + shielding (marginal) \\
\bottomrule
\end{tabular}
\end{center}

\noindent The model plateaued at \(\sim\)15\% between v2.6 and v2.8.
Additional parameters produced negligible improvement, confirming that the
functional form had reached its representational ceiling.

% ====================================================================
\section{Root-Cause Analysis of the 15--17\% Floor}

Four structural causes were identified:

\begin{enumerate}
  \item \textbf{Descriptor degeneracy in the f-block.}
    The descriptor set (cov-radius, period, block, f-electron count, $I_1$)
    maps multiple chemically distinct elements to nearly identical descriptor
    vectors. Irreducible degeneracy of \(\sim\)6--8\% within the f-block.

  \item \textbf{Anomalous elements.}
    Pd ($Z=46$, 4d\textsuperscript{10}), Eu/Am half-shell maxima,
    Yb/No full-shell drops, Lr ($Z=103$, post-actinide spike).
    These elements have electronic configurations that break the smooth
    periodic trends the model assumes.

  \item \textbf{Single-radius basis limitation.}
    The model uses $r_\mathrm{cov}^3$ as its sole volumetric basis.
    Polarizability depends on the outer electron cloud, which is not
    strictly proportional to covalent radius across all blocks.

  \item \textbf{Absence of many-body or correlation physics.}
    Scalar polarizability in heavy atoms is influenced by core-valence
    correlation and spin-orbit effects that no descriptor-free analytical
    model can capture without explicit electronic structure calculation.
\end{enumerate}

% ====================================================================
\section{Decision: Empirical Table as Production Path}

Given the irreducible 15\% floor, a deliberate architectural decision was
made in v2.8.X:

\begin{quote}
\itshape
The empirical table is the primary and sole production path for
$Z = 1$--118.  The generative model is retained only as a fallback for
hypothetical super-heavy elements ($Z > 118$) where no reference data
exists.
\end{quote}

This decision eliminated 100\% of prediction error for all known elements
while preserving the generative path for exploratory use.

% ====================================================================
\section{Anti-Pattern Log}

\begin{center}
\begin{tabular}{lp{7cm}}
\toprule
\textbf{Anti-Pattern} & \textbf{Lesson} \\
\midrule
Over-fitting f-block & Adding more f-block Gaussians beyond half/full-shell
  produced oscillatory fits with worse out-of-sample error. \\
Polynomial Z terms & Replacing $r_\mathrm{cov}^3$ with $Z^n$ polynomials
  destroyed physical interpretability for zero RMS gain. \\
Per-element offsets & Adding per-element additive constants collapsed the
  model to a table with extra steps. \\
Aggressive regularization & L2 penalty strong enough to prevent overfitting
  also prevented the optimizer from finding the 15\% basin. \\
\bottomrule
\end{tabular}
\end{center}

% ====================================================================
\section{SCF Polarization Integration}

The Alpha Model feeds the Drude/SCF polarization solver through:

\begin{equation}
  k_D(Z) = \frac{q_D^2}{\alpha(Z)}
\end{equation}

\noindent where $q_D$ is the Drude particle charge (default 1.0\,\textit{e}).
This couples static polarizability to the iterative self-consistent-field
loop via Thole-style damped dipole interactions.

% ====================================================================
\section{File Index}

\begin{center}
\begin{tabular}{ll}
\toprule
\textbf{File} & \textbf{Purpose} \\
\midrule
\texttt{atomistic/models/alpha\_model.hpp} & Production prediction (table + fallback) \\
\texttt{atomistic/models/atomic\_descriptors.hpp} & Descriptor layer \\
\texttt{atomistic/models/atomic\_descriptors\_ext.hpp} & Extended descriptors \\
\texttt{data/polarizability\_ref.csv} & Reference dataset \\
\texttt{config/alpha\_model\_params.json} & Generative fallback parameters \\
\texttt{tools/fit\_alpha\_model.cpp} & Diagnostic offline fitter \\
\texttt{tools/auto\_fit\_alpha.cpp} & Automated fitting harness \\
\texttt{tests/test\_alpha\_model.cpp} & Invariant tests (T1--T4) \\
\bottomrule
\end{tabular}
\end{center}

% ====================================================================
\section*{References}

\begin{enumerate}
  \item Miller, T.M.\ \textit{JACS} \textbf{112}, 8533 (1990).
  \item Schwerdtfeger, P.; Nagle, J.K.\ \textit{Mol.\ Phys.}\ \textbf{117}, 1200 (2019).
  \item Thole, B.T.\ \textit{Chem.\ Phys.}\ \textbf{59}, 341 (1981).
  \item Lamoureux, G.; Roux, B.\ \textit{J.\ Chem.\ Phys.}\ \textbf{119}, 3025 (2003).
  \item NIST Atomic Spectra Database, \url{https://www.nist.gov/pml/atomic-spectra-database}.
\end{enumerate}

\end{document}
